\subsection*{Chapter 26: Direct Preference Optimization and Its Family}

\paragraph{Exercise 26.1} $\pi^*\propto\piref e^{r/\beta}$. For $\beta=0.1$:
$\pi^*\approx(0.0000,\,0.0001,\,0.9999)$; for $\beta=1$: $(0.179,\,0.292,\,0.529)$; for $\beta=10$:
$(0.465,\,0.308,\,0.227)$. The maximum objective equals $\beta\log Z$ (from the proof of
\cref{prop:dpo:gibbs}): $1.839$, $1.027$, $0.731$ respectively. Sampling random $\pi$ from a Dirichlet
and evaluating $\E_\pi[r]-\beta\KL(\pi\|\piref)$ never exceeds these values, and the gap equals
$\beta\KL(\pi\|\pi^*)$.

\paragraph{Exercise 26.2} BT depends on $r(x,y_w)-r(x,y_l)$, in which $f(x)$ cancels. In
\eqref{eq:dpo:gibbs}, $e^{(r+f)/\beta}=e^{f/\beta}e^{r/\beta}$ and $e^{f(x)/\beta}$ is absorbed by
$Z(x)$, so $\pi^*$ is unchanged. Rewards, and therefore DPO's implicit rewards, are identifiable
only up to a per-prompt offset: compare them only between responses to the same prompt.

\paragraph{Exercise 26.3} With $h=\hat r_w-\hat r_l$, $\partial(-\log\sigma(h))/\partial h=-\sigma(-h)$
and $\nabla_\theta h=\beta(\nabla\log\pi_\theta(y_w\mid x)-\nabla\log\pi_\theta(y_l\mid x))$ (the
reference terms are constants), giving \eqref{eq:dpo:grad}. At initialisation $h=0$ and
$\sigma(0)=1/2$, so $\nabla\mathcal{L}=-\frac{\beta}{2}\nabla(\log\pi_\theta(y_w)-\log\pi_\theta(y_l))$.
This is why the first steps of DPO resemble ``SFT on winners minus SFT on losers'' at half strength.

\paragraph{Exercise 26.4} Setting the derivative $-(1-\varepsilon)\sigma(-h)+\varepsilon\sigma(h)$
to zero gives $\sigma(h^*)=1-\varepsilon$, i.e.\ $h^*=\log\frac{1-\varepsilon}{\varepsilon}$
($\log 9\approx2.20$ for $\varepsilon=0.1$). The minimum is the binary entropy
$-(1-\varepsilon)\log(1-\varepsilon)-\varepsilon\log\varepsilon$. As $\varepsilon\to0$, $h^*\to\infty$:
plain DPO has no finite optimum on a single, always-agreeing pair.

\paragraph{Exercise 26.5} With uniform $\piref$ the reference terms cancel and the log-ratio
margin is $h_0=\log\sigma(z)-\log\sigma(-z)=z$. DPO: $-\log\sigma(\beta z)$ decreases monotonically
in $z$, so the infimum is at $z\to\infty$, i.e.\ $\pi(y_w)\to1$ (deterministic). IPO:
$(z-\frac{1}{2\tau})^2$ is minimised at $z^*=1/(2\tau)$, so $\pi(y_w)=\sigma(1/(2\tau))<1$. The
squared loss keeps the policy at a finite distance from $\piref$, controlled by $\tau$.

\paragraph{Exercise 26.6} Mean log-probs give
$h=\beta\big[\frac{1}{|y_w|}\Delta_w-\frac{1}{|y_l|}\Delta_l\big]$ with
$\Delta=\log\pi_\theta-\log\piref$, i.e.\ DPO with $\beta_y=\beta/|y|$. A winner twice as long as
the loser has its log-ratio down-weighted by $2\times$ relative to the loser's. Compared with
standard DPO, where a long winner accumulates log-ratio over more tokens, this removes the
incentive to lengthen responses. Relative to summed DPO, it favours shorter outputs.

\paragraph{Exercise 26.7} With mismatched truncation, the prompt or response tokens differ between
cache and policy, so at step 0 the per-pair log-ratios are not zero: the loss deviates from
$\ln 2$, the margins are non-zero, and reward accuracy is not $1/2$. Dropout in the policy pass
adds zero-mean noise to the log-ratios, so the loss is slightly \emph{above} $\ln 2$ on average and
accuracy is random around $1/2$. Assertions after the first micro-batch:
\begin{lstlisting}
assert abs(loss0 - math.log(2)) < 1e-3, "policy != reference at step 0"
assert margins0.abs().max() < 1e-4,      "cache/template/truncation mismatch"
assert not policy.training or no_dropout(policy), "disable dropout for DPO"
\end{lstlisting}

\paragraph{Exercise 26.8} Pack $[x, y_w, y_l]$ into one row with positions
$0..T_p{-}1$, $T_p..T_p{+}T_w{-}1$, $T_p..T_p{+}T_l{-}1$, and a mask in which $y_l$ attends to $x$ and
causally to itself but not to $y_w$. The log-probs then match the pair batch to float precision. With
toy responses of about 6 tokens, tokens per pair are $2T_p+12$ (naive) vs $T_p+12$ (shared): 20 vs 16
for $T_p=4$, 44 vs 28 for $T_p=16$, 140 vs 76 for $T_p=64$. The saving approaches 50\% as prompts
dominate.

\paragraph{Exercise 26.9} $2\cdot4\cdot2048\cdot151936\cdot4\approx9.96$\,GB of fp32 logits.
(a) Slicing to the response half before the LM head halves it, to $\approx4.98$\,GB. (b) Chunking over 8
sequence chunks keeps only one chunk's logits alive: $\approx1.24$\,GB peak (plus the saved
per-token log-probs, which are negligible). (c) bf16 logits halve storage, to $\approx4.98$\,GB unchunked;
compute \code{logsumexp} in fp32 per chunk to avoid precision loss.

\paragraph{Exercise 26.10} Per round: $8\times10^4$ responses of $\sim$512 tokens is
$\approx4\times10^7$ generated tokens. Decode for a 7B model is memory-bound
(\cref{ch:inference}); with large batches an H100 sustains on the order of $10^4$ tokens/s, so 8
GPUs take $\approx40\text{M}/(8\cdot1.5\times10^4)\approx5$--6 minutes. Scoring with a 1B reward
model costs $\approx2\cdot10^9\cdot768\cdot8\times10^4\approx1.2\times10^{17}$ FLOPs, under a minute.
Training on $2\times10^4$ pairs of $2(256+512)=1536$ tokens at $8N$ FLOPs/token (policy plus
reference) is $\approx1.7\times10^{18}$ FLOPs, about 9 minutes at 40\% MFU on 8 H100s. The phases
are batched and of similar length, so time-share all 8 GPUs (generate, score, train, sync weights
once per round) rather than splitting them: a disaggregated layout would leave half the GPUs idle
in every phase.
